\documentclass[11pt]{article}
\usepackage[T1]{fontenc}
\usepackage[utf8]{inputenc}
\usepackage{amsmath,amssymb,mathtools}
\usepackage[a4paper,margin=26mm]{geometry}
\usepackage{booktabs,longtable,array}
\usepackage{enumitem}
\usepackage{xcolor}
\usepackage{hyperref}
\allowdisplaybreaks

\hypersetup{
  colorlinks=true,
  linkcolor=blue!55!black,
  citecolor=blue!55!black,
  urlcolor=blue!60!black,
  pdftitle={A Lean transcription of the corrected Lazard quintic formula},
  pdfauthor={ProveIt project}
}

\newcommand{\lean}[1]{\nolinkurl{#1}}
\newcommand{\Q}{\mathbb{Q}}
\newcommand{\K}{\mathbb{K}}
\newcommand{\eps}{\varepsilon}
\newcommand{\ii}[1]{i_{#1}}
\newcommand{\codepath}{\lean{Algebra/PolynomialFormulas/Lean/PolynomialFormulas/LazardQuintic.lean}}
\setlength{\emergencystretch}{3em}

\title{A Lean Transcription of the Corrected Lazard Formula\\
for Solvable Quintics}
\author{ProveIt project}
\date{August 2026}

\begin{document}
\maketitle

\begin{abstract}
Daniel Lazard gave an explicit radical formula for every irreducible solvable
quintic.  This note documents a Lean 4 transcription of the formula, including
Vladimir Reshetnikov's correction of two terms in the printed expression
$P_{22}$ and Valeriy Zapodovnikov's restoration of the exceptional sign choice
needed when the first determination of $\eps$ would make $T=0$.

The Lean development formalizes the conditional data flow after the required
algebraic witnesses have been supplied: normalization and depression,
Lazard's degree-six resolvent, the four equations recovering
$\ii4,\ldots,\ii8$, all auxiliary invariants $D,\ldots,K$, both coherent
sign-selection stages, the corrected $P_{22}$, and inverse Fourier
reconstruction of five candidate outputs.  It proves the normalization identity, the
equivalence of the monic and fraction-free resolvent equations, preservation of
$\eps^2$ by the correction, nonvanishing consequences of the branch
certificates, coherence of all four quadratic sign branches, and soundness of
finite branch selection.  It does \emph{not} prove that the five displayed
outputs are roots, nor that the needed certificates exist.  Pointwise
soundness and full five-factor completeness are recorded as separate
propositions, not as axioms.
\end{abstract}

\tableofcontents

\section{Purpose, sources, and claim boundary}

The formal source is
\begin{center}
  \codepath.
\end{center}
It is imported by the public \lean{PolynomialFormulas} module.  The source
materials were reviewed at Smithereens commit
\lean{da7026ec1369ade49bb6d964019a4c124fdd5091}.  The principal references are:

\begin{enumerate}[label=(\arabic*)]
\item Lazard's chapter \emph{Solving Quintics by Radicals}, especially the
  explicit formula on printed pp.~219--222 \cite{Lazard2004};
\item Lazard's original Maple worksheet, version 3 (March 2002), distributed
  with those materials;
\item Reshetnikov's Wolfram-language transcription and warning about the two
  printed $P_{22}$ errors \cite{MSE};
\item Zapodovnikov's corrected 2025 version on the same page, especially
  \lean{If[(g+h/F)==0,F=-F]}; and
\item Dummit's independent solvability resolvent \cite{Dummit1991}, together
  with Chapman's repair of the specialized-distinctness step
  \cite{ChapmanNote}.
\end{enumerate}

The expanded review also covered every text-readable document under
\lean{src/Polynomial} whose path or contents mention ``quintic'', including the
Reshetnikov notebooks, the Wolfram notebooks, the independent
\lean{solving-quintics} package, Dummit's paper, and the corpus reports.  These
additional files provide useful examples and independent checks, but they do
not supersede Lazard's formula.  In particular, the
\lean{solving-quintics} package switches to numerical root matching after the
resolvent and is not an exact implementation of the formula documented here.

The present claim is deliberately limited:

\begin{quote}
The Lean definitions are an exact, type-correct transcription of the corrected
algebraic algorithm, organized so that its eventual correctness proof can be
stated without changing the API.  No theorem in this contribution asserts the
unproved final root identity.
\end{quote}

This distinction matters.  Lean accepts division by zero as a total field
operation, and so a raw expression alone would not faithfully express
Lazard's partial algebraic construction.  The formalization therefore records
every required denominator condition and every shared radical determination in
\lean{RadicalCertificate}.

\section{Mathematical scope}

Lazard begins with an irreducible quintic over a field of characteristic
different from $2$ and $5$.  His compact expanded formula uses a particular
fourth projection whose invertibility is justified when $-1$ is not a square;
this is harmless over $\Q$.  The paper also sketches an alternate projection
for more general fields, but does not print its expanded formula.

The Lean expressions are polymorphic over a field $K$ with
\lean{CharZero K}.  This is a slightly stronger characteristic assumption and
ensures that all displayed rational denominators are meaningful.  The actual
rational front end is made explicit by
\lean{RationalInvariantCertificate}: the tuple
$(\ii4,\ldots,\ii8)$ is first certified in $\Q$ and only then mapped to a
radical extension.  The abstract generic API remains useful for algebraic
lemmas.

The generic propositions \lean{FormulaSound} and \lean{FormulaFactors} are
deliberately stronger \emph{conditional algebraic targets}: they quantify over
an arbitrary characteristic-zero field, but only after every displayed
invariant and radical certificate has been supplied.  They do not assert
existence of those certificates.  Thus Lazard's sourced irreducibility and
$-1$ nonsquare hypotheses remain obligations for the applicability/existence
bridge, while the rational specialization fixes the coefficient field to
$\Q$ and performs radical calculations in an extension field.

For a complete solver on arbitrary quintics, reducible inputs must be factored
first.  Every proper factor has degree at most four and can in principle be
handled by the existing low-degree formulas.  This module implements neither
that factor search nor the rational-root search for the sextic: its rational
entry point accepts a supplied \lean{RationalInvariantCertificate}.  Lazard's
core is the irreducible branch; the formalization does not misrepresent it as
a total reducer for arbitrary degree-five polynomials.

\section{The algorithm at a glance}

For an irreducible rational quintic, the corrected construction is:

\begin{enumerate}
\item normalize and translate to $Y^5+pY^3+qY^2+rY+s$;
\item find the rational root $\ii4$ of Lazard's sextic resolvent;
\item solve the four linear equations of Figure~3 for
  $\ii5,\ii6,\ii7,\ii8$;
\item calculate $D,E,F,G$ and choose $\eps^2=5D$;
\item if $E+F/\eps=0$, replace $\eps$ by $-\eps$;
\item choose $T^2=\frac52(E+F/\eps)$ and set $U=5G/(T\eps)$;
\item among four coherent sign triples $(\eps,T,U)$, choose one for which
  $Q_1\ne0$;
\item choose one shared fifth root $P_1^5=Q_1$, calculate the corrected
  $P_2,P_3,P_4$, and apply the inverse discrete Fourier transform;
\item translate every resulting root back by $-b/(5a)$.
\end{enumerate}

The group-theoretic tower behind these steps is
\[
 C_5\subset D_5\subset F_{20}\subset S_5,
 \qquad [D_5:C_5]=2,\quad [F_{20}:D_5]=2,\quad [S_5:F_{20}]=6.
\]
The sextic selects a conjugate of the maximal solvable subgroup $F_{20}$;
the two square roots descend through the two index-two inclusions; one fifth
root resolves the remaining cyclic action.

\section{Normalization and depression}

Write
\[
 f_0(X)=aX^5+bX^4+cX^3+dX^2+eX+f,\qquad a\ne0.
\]
Under $X=Y-b/(5a)$ and division by $a$, one obtains
\[
 f(Y)=Y^5+pY^3+qY^2+rY+s,
\]
where
\begin{align*}
p&=\frac{5ac-2b^2}{5a^2},\\
q&=\frac{25a^2d-15abc+4b^3}{25a^3},\\
r&=\frac{125a^3e-50a^2bd+15ab^2c-3b^4}{125a^4},\\
s&=\frac{3125a^4f-625a^3be+125a^2b^2d-25ab^3c+4b^5}
        {3125a^5}.
\end{align*}

These are \lean{depress}.  Lean proves, for every $Y$,
\[
 f_0\!\left(Y-\frac{b}{5a}\right)=a\,f(Y)
\]
as \lean{depress_eval}.  It also proves that depression commutes with a field
homomorphism (\lean{depress_map}), which is the bridge from rational
certificates to the radical field.

\section{Fourier invariants and metacyclic data}

Let $x_0,\ldots,x_4$ be an ordering of the roots and let $\omega$ be a
primitive fifth root of unity.  Lazard uses the Fourier sums
\[
 s_k=\sum_{j=0}^{4}\omega^{kj}x_j,\qquad k=0,\ldots,4.
\]
The depressed condition gives $s_0=0$.  Define cyclic invariants
\[
 S_1=s_1^5,\qquad S_2=s_2s_1^3,\qquad
 S_3=s_3s_1^2,\qquad S_4=s_4s_1.
\]
Once $P_1^5=S_1$ is chosen with $P_1\ne0$, the remaining Fourier components
are $P_2=S_2/P_1^3$, $P_3=S_3/P_1^2$, and $P_4=S_4/P_1$.

The rational metacyclic basis used to express these quantities is
\begin{align*}
\ii4&=\sum_C x_0^2(x_1x_4+x_2x_3),&
\ii5&=\sum_C x_0^3(x_1x_4+x_2x_3),\\
\ii6&=\sum_C x_0^4(x_1x_4+x_2x_3),&
\ii7&=\sum_C x_0^3(x_1^2x_4^2+x_2^2x_3^2),\\
\ii8&=\sum_C x_0^4(x_1^2x_4^2+x_2^2x_3^2).&&
\end{align*}
Lean stores their values in \lean{Invariants}.  Their origin in a root tuple
will be established only in a later correctness bridge; this formula module
uses the exact resolvent and linear relations as its semantic interface.

\section{The sextic resolvent}

The discriminant of $X^5+pX^3+qX^2+rX+s$ is
\begin{align*}
\Delta={}&108p^5s^2-72p^4qrs+16p^4r^3+16p^3q^3s-4p^3q^2r^2
 -900p^3rs^2\\
&+825p^2q^2s^2+560p^2qr^2s-128p^2r^4-630pq^3rs
 +144pq^2r^3\\
&-3750pqs^3+2000pr^2s^2+108q^5s-27q^4r^2
 +2250q^2rs^2\\
&-1600qr^3s+256r^5+3125s^4.
\end{align*}
This is \lean{discriminant}; it was independently compared with a computer
algebra discriminant on all $625$ tuples $p,q,r,s\in\{-2,-1,0,1,2\}$ during
the source audit.

For a variable $z=\ii4$, put
\begin{align*}
A(z)={}&2z^3+8rz^2+(-6p^2r+2pq^2-50qs+24r^2)z\\
&-15p^2qs-16p^2r^2+13pq^2r+125ps^2
-2q^4-200qrs+64r^3.
\end{align*}
Then Lazard's monic sextic is
\begin{equation}\label{eq:resolvent}
 R(z)=\left(\frac{A(z)}2\right)^2-
 \left(z+3r+\frac{p^2}{4}\right)\Delta.
\end{equation}
The fraction-free equation used by the corrected Wolfram code is
\begin{equation}\label{eq:resolvent-ff}
 A(z)^2=(4z+p^2+12r)\Delta.
\end{equation}
Lean's two definitions are \lean{resolventEval} and
\lean{resolventPolynomial}.  Their agreement under evaluation is the theorem
\lean{resolventPolynomial_eval}.  The theorem
\lean{resolventEval_eq_zero_iff} proves the equivalence of
\eqref{eq:resolvent} and \eqref{eq:resolvent-ff}.

There is a notation trap in the printed sources.  Cayley's earlier invariant is
\[
 \Theta=(V-W)^2=4\ii4+p^2+12r,
\]
whereas the section-7 sextic variable is $z=\ii4$.  Lean uses separate names
\lean{cayleyTheta} and \lean{resolventCore}, preventing this collision.

For an irreducible solvable rational quintic, the sextic has a rational root.
Dummit's independent $F_{20}$ resolvent is the same invariant and gives the
same criterion.  A future proof must not infer specialized distinctness from
generic transitivity alone: Chapman's note supplies the missing argument, and
Lazard gives a separability proof for the Cayley resolvent.

\section[Recovering i5 through i8]{Recovering $\ii5,\ldots,\ii8$}

After choosing $\ii4$, Lazard's Figure~3 gives four equations linear in
$\ii5,\ldots,\ii8$.  They are transcribed as
\lean{i4SquareRhs}, \lean{i4CubeRhs}, \lean{i4FourthRhs}, and
\lean{i4FifthRhs}; together with the sextic equation they form
\lean{InvariantRelations}.

\begin{align*}
\ii4^2={}&5\ii8-2p\ii6+4q\ii5-2p^2\ii4-6p^2r+2pq^2+10qs+4r^2,\\
\ii4^3={}&\frac12\bigl[(3p^2-20r)\ii8+(-pq-50s)\ii7
 +(-3p^3+28pr-12q^2)\ii6\\
& +(3p^2q-45ps-6qr)\ii5
 +(-3p^4+36p^2r-15pq^2+60qs-32r^2)\ii4\\
& -6p^4r+3p^3q^2+41p^2qs+52p^2r^2-54pq^2r-250ps^2
 +14q^4+140qrs-80r^3\bigr],\\
\ii4^4={}&(19p^2r-9pq^2+225qs-60r^2)\ii8
 +(15p^2s-8pqr+3q^3+100rs)\ii7\\
&+(-4p^3r+4p^2q^2-105pqs-16pr^2+29q^2r+125s^2)\ii6\\
&+(-9p^3s+17p^2qr-8pq^3+140prs+155q^2s-68qr^2)\ii5\\
&+(-4p^4r+4p^3q^2-79p^2qs-16p^2r^2+15pq^2r-25ps^2
  +4q^4+80qrs)\ii4\\
&+6p^4qs-22p^4r^2+16p^3q^2r-4p^2q^4-404p^2qrs+68p^2r^3\\
&+132pq^3s+42pq^2r^2+550prs^2-30q^4r-50q^2s^2+20qr^2s+16r^4,
\end{align*}
and
{\small
\begin{align*}
2\ii4^5={}&
(15p^4r-5p^3q^2+290p^2qs-152p^2r^2-27pq^2r-1375ps^2
 +22q^4-700qrs+240r^3)\ii8\\
&+(18p^4s-11p^3qr+3p^2q^3-530p^2rs+110pq^2s+124pqr^2
 -41q^3r-2375qs^2+200r^2s)\ii7\\
&+(-15p^5r+5p^4q^2-212p^3qs+168p^3r^2-83p^2q^2r
 +325p^2s^2+10pq^4+1560pqrs\\
&\qquad-176pr^3-620q^3s-12q^2r^2-1500rs^2)\ii6\\
&+(15p^4qr-5p^3q^3-147p^3rs+351p^2q^2s-90p^2qr^2
 -43pq^3r-3175pqs^2\\
&\qquad-420pr^2s+20q^5+215q^2rs+152qr^3+625s^3)\ii5\\
&+(-15p^6r+5p^5q^2-200p^4qs+200p^4r^2-110p^3q^2r
 +355p^3s^2+15p^2q^4\\
&\qquad+1728p^2qrs-432p^2r^3-752pq^3s+220pq^2r^2-200prs^2
 -43q^4r+1825q^2s^2\\
&\qquad-2640qr^2s+512r^4)\ii4\\
&-30p^6r^2+25p^5q^2r+198p^5s^2-5p^4q^4-491p^4qrs+364p^4r^3
 +181p^3q^3s\\
&-286p^3q^2r^2-810p^3rs^2+95p^2q^4r+3005p^2q^2s^2
 +4120p^2qr^2s-1088p^2r^4\\
&-12pq^6-4095pq^3rs+612pq^2r^3-15875pqs^3+900pr^2s^2
 +858q^5s-34q^4r^2\\
&+10700q^2rs^2-6240qr^3s+960r^5\\
&\qquad+6250s^4.
\end{align*}
}

Retaining these as equations, instead of embedding the enormous pre-solved
rational functions, is intentional.  It is exactly the finite linear-system
interface needed for a later proof of existence and uniqueness.

\section{The two quadratic stages}

The explicit metacyclic invariants on Lazard p.~221 are:
\begin{align*}
D={}&40p\ii8-120q\ii7+(-24p^2+100r)\ii6+(88pq-300s)\ii5\\
&+(-24p^3+100pr+24q^2)\ii4-80p^3r+40p^2q^2-480pqs
 +160pr^2+332q^2r+125s^2,\\
E={}&(3p^2+20r)\ii6+(-pq-50s)\ii5+(3p^3+12pr+3q^2)\ii4\\
&+4p^3r-3p^2q^2+40pqs+16pr^2-21q^2r+125s^2,\\
F={}&(-65p^2q+875ps-550qr)\ii8+(-58p^2r+41pq^2-275qs+440r^2)\ii7\\
&+(85p^3q-520p^2s-298pqr+366q^3+2100rs)\ii6\\
&+(4p^3r-73p^2q^2+2095pqs-56pr^2-748q^2r-4875s^2)\ii5\\
&+(85p^4q-418p^3s-440p^2qr+419pq^3+1590prs-1040q^2s+524qr^2)\ii4\\
&-12p^5s+158p^4qr-85p^3q^3-1462p^3rs-159p^2q^2s
 +142p^2qr^2\\
&+896pq^3r+175pqs^2+2900pr^2s-402q^5-1925q^2rs-448qr^3-1875s^3,\\
G={}&(-35p^2q-250ps-200qr)\ii8+(-22p^2r+19pq^2+650qs-40r^2)\ii7\\
&+(15p^3q+195p^2s+68pqr-6q^3-1100rs)\ii6\\
&+(-4p^3r-27p^2q^2-270pqs+96pr^2-182q^2r+3000s^2)\ii5\\
&+(15p^4q+213p^3s+50p^2qr+pq^3-940prs+515q^2s-184qr^2)\ii4\\
&+12p^5s+42p^4qr-15p^3q^3+492p^3rs-156p^2q^2s+358p^2qr^2\\
&-246pq^3r+2825pqs^2-1400pr^2s+42q^5+550q^2rs-232qr^3-1250s^3.
\end{align*}
They are \lean{invariantD}, \lean{invariantE}, \lean{invariantF}, and
\lean{invariantG}.

Choose $\eps_0$ with $\eps_0^2=5D$.  The corrected choice is
\begin{equation}\label{eq:eps-correct}
 \eps=\begin{cases}
 -\eps_0,&E+F/\eps_0=0,\\
 \eps_0,&E+F/\eps_0\ne0.
 \end{cases}
\end{equation}
\begin{samepage}
This is literally \lean{correctEpsilon}.  Lean proves
\lean{correctEpsilon_sq}.  Under the sourced assertion that at least one of
$E+F/\eps_0$ and $E-F/\eps_0$ is nonzero, the theorem
\begin{center}
\lean{correctedTRadicand_ne_zero}
\end{center}
proves that the corrected radicand is nonzero.  One then chooses
\begin{equation}\label{eq:TU}
 T^2=\frac52\left(E+\frac F\eps\right),\qquad
 U=\frac{5G}{T\eps}.
\end{equation}
\end{samepage}
The factor $5/2$ is the one used in section~7, Maple, and both Wolfram ports;
an earlier prose passage in Lazard omits the relevant scaling.

The records \lean{EpsilonChoice}, \lean{TChoice}, and
\lean{RadicalCertificate} expose the equations and nondegeneracy assumptions.
In particular, the certificate stores the complementary equation
$U^2=\frac52(E-F/\eps)$.  The theorem
\lean{RadicalCertificate.t_nonzero} derives $T\ne0$ from the corrected
radicand rather than trusting total field division, and
\lean{RadicalCertificate.initial_relations} proves $TU\eps=5G$ from the
definition of $U$.

\section{The fifth-root radicand and coherent signs}

Define
\begin{align*}
H={}&25(2\ii5-pq-5s),\\
I={}&25\bigl[40p\ii8-70q\ii7+(-24p^2+100r)\ii6+(68pq-300s)\ii5\\
&\qquad+(-24p^3+100pr-46q^2)\ii4-80p^3r+20p^2q^2-255pqs
 +160pr^2-28q^2r+125s^2\bigr],\\
J={}&-25p\ii8-25q\ii7+(-9p^2-60r)\ii6+(-7pq+525s)\ii5\\
&+(-p^3-96pr+11q^2)\ii4+50p^3r-7p^2q^2-145pqs-308pr^2
 +128q^2r-1000s^2,\\
K={}&-125p\ii8+75q\ii7+(67p^2-420r)\ii6+(-109pq+1175s)\ii5\\
&+(63p^3-412pr+27q^2)\ii4+210p^3r-79p^2q^2-415pqs-676pr^2
 +496q^2r-750s^2.
\end{align*}
These are \lean{invariantH} through \lean{invariantK}.  For a coherent
triple $(\eps,T,U)$, set
\begin{equation}\label{eq:q1}
 Q_1=\frac54\left(H+\frac I\eps+\frac{TJ+UK}{E}\right).
\end{equation}

The four allowed triples are
\[
(\eps,T,U),\quad(\eps,-T,-U),\quad(-\eps,U,-T),\quad(-\eps,-U,T).
\]
They are the four constructors handled by \lean{branchTriple}.  The selector
\lean{selectSignBranch} returns the first branch with $Q_1\ne0$, or
\lean{none} if the mathematical existence hypothesis has not been supplied.
Its theorem \lean{selectSignBranch_sound} proves that a returned branch is
indeed nonzero.  Quadratic coherence is the separate theorem
\begin{center}
\lean{QuadraticRelations.of_branchTriple}.
\end{center}
It proves that each transformation preserves both square equations and
$TU\eps=5G$.  A \lean{RadicalCertificate} stores the chosen branch, the
nonzero fact, and one shared $P_1$ satisfying $P_1^5=Q_1$; Lean derives
$P_1\ne0$ in \lean{RadicalCertificate.p1_nonzero}.

This finite branch mechanism is the algebraic content of the corrected
algorithm.  Domen's later use of Mathematica's \lean{Together} merely repairs
syntactic evaluation of fractional powers during root matching.  Since Lean
uses a single $\omega$ and field expressions, it has no mathematical analogue
here.

\section[The corrected P22 and Fourier reconstruction]
  {The corrected $P_{22}$ and Fourier reconstruction}

Lazard's remaining numerators are transcribed one per Lean definition:
\begin{align*}
P_{41}&=-5p,\\
P_{42}&=5(10\ii7-4p\ii5-14q\ii4-4p^2q+45ps-72qr),\\
P_{31}&=-25q,\\
P_{32}&=25(-10\ii8+2p\ii6-22q\ii5+2p^2\ii4+20p^2r+2pq^2-35qs-40r^2),\\
P_{33}&=5(35\ii8-4p\ii6+23q\ii5+(-6p^2+12r)\ii4
 -58p^2r+14pq^2-105qs+76r^2),\\
P_{34}&=5(5\ii8-22p\ii6+14q\ii5+(-18p^2+16r)\ii4
 -34p^2r+22pq^2-140qs+68r^2),\\
P_{21}&=5(3\ii4+2p^2-16r),\\
P_{22}&=25\bigl(-10q\ii6+(8p^2-50r)\ii5+(-2pq-25s)\ii4\\
&\hspace{35mm}+\boxed{8p^3q}+\boxed{70q^3}-20p^2s-26pqr+50rs\bigr),\\
P_{23}&=25(-4p\ii7-q\ii6+4r\ii5+(-3pq+15s)\ii4
 +26p^2s-26pqr+7q^3-40rs),\\
P_{24}&=25(3p\ii7-18q\ii6+22r\ii5+(-14pq+20s)\ii4
 +18p^2s-33pqr+21q^3+30rs).
\end{align*}

The boxed terms are Reshetnikov's corrections.  Lazard p.~222 prints $8p^3$
and $70q^3q$; the corrected terms are independently confirmed by the original
Maple worksheet, which has $200p^3q$ and $1750q^3$ after distributing the
outer factor $25$.  Lean makes the corrected formula conspicuous in
\lean{p22}.

With the selected $(\eps,T,U)$ and shared $P_1$, define
\begin{align*}
P_4={}&\frac{P_{41}}{2P_1}+\frac{P_{42}}{2\eps P_1},\\
P_3={}&\frac{P_{31}}{4P_1^2}+\frac{P_{32}}{4\eps P_1^2}
 +\frac{P_{33}T+P_{34}U}{10EP_1^2},\\
P_2={}&\frac{P_{21}}{4P_1^3}+\frac{P_{22}}{4\eps P_1^3}
 +\frac{P_{23}T+P_{24}U}{10EP_1^3}.
\end{align*}
These are \lean{fourierP4}, \lean{fourierP3}, and \lean{fourierP2}.

For one primitive fifth root $\omega$, all five depressed candidates are
\begin{equation}\label{eq:dft}
 y_k=\frac{\omega^kP_1+\omega^{2k}P_2+
                   \omega^{3k}P_3+\omega^{4k}P_4}{5},
 \qquad k=0,1,2,3,4.
\end{equation}
The original roots are
\[
 x_k=y_k-\frac{b}{5a}.
\]
Lean represents $\omega$ by \lean{FifthRootOfUnity}; its field
\lean{primitive} certifies that \lean{value} is a primitive fifth root.  The
single indexed formula \lean{solveDepressed} replaces five branch-sensitive
Mathematica expressions.  \lean{solveGeneral} performs the final translation,
and \lean{solveRational} keeps the rational invariant certificate visible.

\section{Lean-to-paper crosswalk}

\small
\begin{longtable}{>{\raggedright\arraybackslash}p{0.34\textwidth}
                  >{\raggedright\arraybackslash}p{0.22\textwidth}
                  >{\raggedright\arraybackslash}p{0.35\textwidth}}
\toprule
Lean declaration & Lazard notation/source & Role \\
\midrule
\endfirsthead
\toprule
Lean declaration & Lazard notation/source & Role \\
\midrule
\endhead
\lean{GeneralQuintic}, \lean{GeneralQuintic.polynomial} & $a,b,c,d,e,f$ & Original polynomial and its coefficient representation.\\
\lean{GeneralQuintic.polynomial_eval} & evaluation at $x$ & Proved agreement of polynomial and scalar evaluation.\\
\lean{DepressedQuintic}, \lean{depress} & $p,q,r,s$; p.~219 & Normalize and depress.\\
\lean{depress_eval} & $X=Y-b/(5a)$ & Kernel-checked translation identity.\\
\lean{discriminant} & $\Delta$; pp.~219--220 & Explicit depressed-quintic discriminant.\\
\lean{resolventCore} & $A(z)$ & Cubic whose square enters the sextic.\\
\lean{cayleyTheta} & $\Theta=4\ii4+p^2+12r$; p.~219 & Keeps Cayley $\Theta$ distinct from $z=\ii4$.\\
\lean{resolventPolynomial} & $R(z)$; pp.~219--220 & Monic degree-six polynomial.\\
\lean{resolventEval_eq_zero_iff} & fraction-free Stack Exchange equation & Proved equivalence of the two displays.\\
\lean{Invariants} & $\ii4,\ldots,\ii8$; pp.~217--220 & Metacyclic invariant values.\\
\lean{i4SquareRhs}--\lean{i4FifthRhs} & Figure~3; pp.~220--221 & Linear system for $\ii5,\ldots,\ii8$.\\
\lean{InvariantRelations}, \lean{InvariantRelations.map} & sextic plus Figure~3 & Exact rational-stage certificate, proved stable under scalar extension.\\
\lean{invariantD}--\lean{invariantG} & $D,E,F,G$; p.~221 & Quadratic-stage invariant polynomials.\\
\lean{correctEpsilon} & $\eps\mapsto-\eps$ if $T=0$; p.~221 and Zapodovnikov & Corrected first sign choice.\\
\lean{squareBranchAvailable} & not both $E\pm F/\eps_0$ zero & Explicit nondegeneracy obligation.\\
\lean{correctedTRadicand_ne_zero} & p.~221 nonvanishing argument & Proved consequence of the obligation.\\
\lean{radicalU} & $U=5G/(T\eps)$ & Second quadratic-stage value.\\
\lean{invariantH}--\lean{invariantK} & $H,I,J,K$; pp.~221--222 & Projections defining $Q_1$.\\
\lean{QuadraticRelations} & equations for $\eps,T,U$ & Records both squares and $TU\eps=5G$.\\
\lean{SignBranch}, \lean{branchTriple} & four coherent triples & Finite shared-sign orbit.\\
\lean{QuadraticRelations.of_branchTriple} & coherent sign action & Proves all four transformations preserve the equations.\\
\lean{selectSignBranch} & paper's $Q_1\ne0$ choice & Honest \lean{Option}; no unsafe list indexing.\\
\lean{q1} & $Q_1$ & Fifth-root radicand.\\
\lean{p22} & corrected $P_{22}$; p.~222/MSE & Contains $8p^3q$ and $70q^3$.\\
\lean{fourierP2}--\lean{fourierP4} & $P_2,P_3,P_4$; p.~222 & Remaining Fourier components.\\
\lean{RadicalCertificate} & shared $\eps,T,U,P_1$ & Equations and nonzero obligations stored once.\\
\lean{RadicalCertificate.initial_relations}, \lean{RadicalCertificate.chosen_relations} & quadratic coherence & Connect the certificate to both the initial and selected triples.\\
\lean{FifthRootOfUnity} & primitive $\omega$ & Shared fifth root of unity.\\
\lean{solveDepressed} & inverse Fourier formula & All five depressed candidates.\\
\lean{solveGeneral} & $x_k=y_k-b/(5a)$ & Undo translation.\\
\lean{RationalInvariantCertificate} & rational $\ii4,\ldots,\ii8$ & A supplied witness over $\Q$; it is not a search procedure.\\
\lean{FormulaSound}, \lean{RationalFormulaSound} & pointwise root identity & Unproved propositions; not axioms and not asserted.\\
\lean{FormulaFactors}, \lean{RationalFormulaFactors} & full five-factor identity & Stronger unproved completeness targets.\\
\bottomrule
\end{longtable}
\normalsize

\section{Why the certificates are part of the algorithm}

The following denominators occur in the formula:
\[
 a,\qquad \eps,\qquad T,\qquad E,\qquad P_1.
\]
The formalization treats their nonvanishing as follows.

\begin{itemize}
\item $a\ne0$ is an explicit hypothesis of \lean{depress_eval} and of the
  future soundness and factorization targets.
\item The field \lean{epsilon0_nonzero} records the sourced irreducible-case
  fact; the theorem \lean{epsilon_nonzero} in the
  \lean{RadicalCertificate} namespace proves that sign correction preserves it.
\item \lean{epsilon_branches}, the square equation for $T$, and
  \lean{correctedTRadicand_ne_zero} imply $T\ne0$.
\item The stored square equation for $U$, together with the definition
  $U=5G/(T\eps)$, yields \lean{RadicalCertificate.initial_relations};
  \lean{QuadraticRelations.of_branchTriple} proves that all four coherent
  sign choices preserve the two square equations and $TU\eps=5G$.
\item \lean{invariantE_nonzero} records Lazard's $E=T^2+U^2\ne0$ argument in
  the rational case.
\item \lean{branch_nonzero} and $P_1^5=Q_1$ imply $P_1\ne0$.
\end{itemize}

This organization is intentionally relational.  A later executable algebraic-
number layer may compute the certificates, but the mathematical formula is not
made dependent on a particular computer algebra system's analytic square-root
or fractional-power convention.

\section{What remains to prove}

The generic target \lean{FormulaSound} says that, for a valid leading
coefficient, an invariant tuple satisfying \lean{InvariantRelations}, a valid
radical certificate, and a primitive fifth root of unity, every indexed value
from \lean{solveGeneral} annihilates the original quintic.  Its rational
specialization \lean{RationalFormulaSound} additionally requires that the
invariant tuple be supplied and certified in $\Q$ before scalar extension.
Neither proposition asserts that such witnesses can be found.

Pointwise soundness alone does not express that all five roots, with their
multiplicities, have been recovered.  The separate targets
\lean{FormulaFactors} and \lean{RationalFormulaFactors} state the stronger
polynomial identity
\[
 \text{\lean{GeneralQuintic.polynomial}}
   =a\prod_{k=0}^{4}\bigl(X-\text{\lean{solveGeneral}}(k)\bigr).
\]
Thus the API cleanly separates conditional root soundness, five-root
completeness, and the still-absent existence/search layer.

A complete proof can be decomposed into the following independent lemmas:

\begin{enumerate}
\item identify the sextic with the $F_{20}$ orbit resolvent and prove the
  rational-root solvability criterion;
\item prove specialized separability/distinctness, using Lazard's Cayley
  argument or Chapman's repair rather than the invalid transitivity shortcut;
\item derive the Figure~3 linear system from the root-defined invariants and
  prove its determinant nonzero;
\item derive the displayed $D,\ldots,K$ reductions;
\item prove that the two quadratic stages recover the cyclic projections;
\item prove that one of the four coherent $Q_1$ values is nonzero;
\item identify $P_1,\ldots,P_4$ with the four Fourier sums; and
\item apply Fourier inversion and \lean{depress_eval}.
\end{enumerate}

The factorization target then gives exhaustiveness and, under irreducibility
and separability, pairwise distinctness.  None of these unfinished statements
is admitted with \lean{axiom},
\lean{sorry}, or an opaque external oracle.

\section{Regression evidence and source cautions}

The wider Smithereens review found useful future regression cases, but also
showed why notebook output is not a proof oracle.

\begin{itemize}
\item The corrected zero-$T$ branch is exercised by
  $X^5-10X^2-10X-6$, for which the audited source gives $\ii4=30$.
\item An older direct notebook records a failure on $X^5-5X-12$ when it
  omits an adequate branch fallback.  This is a valuable all-five-roots case.
\item Reshetnikov's strongest stored exact harness checks five values,
  substitution residuals, and distinctness for six polynomials, including the
  very large target polynomial from the Stack Exchange question.
\item Some other notebook runs abort after ten cases, some use
  \lean{PossibleZeroQ}, and the independent package checks high-precision
  residuals without proving symbolic identities or complete root coverage.
\item The label of $X^5+20X+32$ as $F_{20}$ in one newer document is
  incorrect; its group is dihedral.  It should not be used as an $F_{20}$
  classification fixture.
\item Dummit's printed examples have later corrections, and his original
  factorization proof needs Chapman's specialized-distinctness lemma.
\end{itemize}

Consequently the committed Lean module uses source formulas and algebraic
certificates, not notebook success labels.  Future regression computations are
useful checks against transcription errors, but kernel-checked symbolic
identities will remain the correctness criterion.

\section{Reproducibility and audit}

From the repository root, the module and its public import surface are checked
with
\begin{verbatim}
lake env lean \
  Algebra/PolynomialFormulas/Lean/PolynomialFormulas/LazardQuintic.lean

LEAN_NUM_THREADS=0 lake build +PolynomialFormulas
\end{verbatim}
The document is rendered by first changing to its directory and then running
the following command three times (for the table of contents, citations, and
cross-references):
\begin{verbatim}
cd Algebra/PolynomialFormulas
pdflatex -interaction=nonstopmode -halt-on-error \
  LazardQuinticFormalization.tex
\end{verbatim}

For the present contribution, an audit should additionally verify:
\begin{enumerate}
\item no occurrence of \lean{sorry}, \lean{admit}, or a new \lean{axiom};
\item the displayed corrected terms occur as
  \lean{8 * c.p ^ 3 * c.q} and \lean{70 * c.q ^ 3};
\item \lean{correctEpsilon} precedes construction of $T$ and $U$;
\item \lean{solveDepressed} uses one stored $P_1$ and one primitive $\omega$;
  and
\item the PDF is built from the committed \LaTeX{} source without errors.
\end{enumerate}

\begin{thebibliography}{9}

\bibitem{Lazard2004}
Daniel Lazard,
\emph{Solving Quintics by Radicals},
in Olav Arnfinn Laudal and Ragni Piene (eds.),
\emph{The Legacy of Niels Henrik Abel}, Springer, 2004, pp.~207--225.
\url{https://doi.org/10.1007/978-3-642-18908-1_6}.

\bibitem{Dummit1991}
David S. Dummit,
\emph{Solving Solvable Quintics},
Mathematics of Computation 57 (1991), no.~195, 387--401.

\bibitem{ChapmanNote}
Robin Chapman,
\emph{Note on Solving Solvable Quintics},
one-page correction note hosted with Dummit's source materials.
\url{https://site.uvm.edu/ddummit/files/2021/04/Solving_Solvable_Quintics_note_by_Robin_Chapman.pdf}.

\bibitem{MSE}
Vladimir Reshetnikov and Valeriy Zapodovnikov,
\emph{Solving quintic in radicals},
Mathematica Stack Exchange question 133550 and answers, 2016--2025.
\url{https://mathematica.stackexchange.com/questions/133550/solving-quintic-in-radicals}.

\end{thebibliography}

\end{document}
